\section{Equal weights: the gap survives}\label{sec:tone}

This section is about the one operator that \cref{sec:closing} leaves open,
\[
  T=J_1\cdot(J_2\times J_3+J_4\times J_5),\qquad A=-8iT,\qquad q=8|\text{eigenvalue of }T|,
\]
on $\Inv(j_1\otimes\cdots\otimes j_5)$, with $s=j_1$ the distinguished spin.

The chain of \cref{fig:ssh} already says why equal weights should keep the gap: the
inner pair's links are the weak ones, so the dimerisation is trivial and nothing
localises at the ends. For $j_1=\frac12$ that picture is the whole story. For larger
$j_1$ the chain widens into the strip of \cref{fig:grid}, and the question is whether
anything in two dimensions can do what the topological dimerisation did in one.

The answer we can prove is no on every class listed below, and the plan is the same
each time. Up to a positive rescaling, the strip carries two Stark matrices, whose
spectra are integers (\cref{subsec:heronstark}). Where the triangle inequalities do not
cut the strip, that settles the matter (\cref{subsec:closed}). Where they do, we look,
for a suitable pairing of the legs, for a one-dimensional object (a chain, a
two-point pencil, a pair of coupled chains, a tridiagonal stage) whose explicit inverse
is a product of link ratios, and show that the ratios do not compound. That is the
quantitative form of trivial dimerisation. What changes from slice to slice is only
the shape of the one-dimensional object: for $j_1=\frac12$ the strip is a chain; for
$j_1=1$ one colour class sits on a single diagonal and the problem becomes a pencil;
for half-integer $j_1$ the diagonals pair up into zigzag chains coupled through
diagonal matrices; for integer $j_1\ge2$ one diagonal is left over and is shared. The
difficulty grows with $j_1$ because the number of coupled chains does, and the losses
of each coupling compound: at $j_1=\frac32$ the bound left over is close enough to
the threshold that a finite range of parameters has to be checked by computer, and at
$j_1=\frac92$ the simple chained inequality no longer certifies every tuple.

\begin{theorem}{The gap at equal weights}{tone}
Every nonzero eigenvalue of $8T$ satisfies $q\ge q_*$ in each of the following cases.
\begin{enumerate}
\item[(a)] $j_1=\frac12$, with $q_*=6\sqrt2/\pi=2.7009$ \textup{(}\cref{thm:B2}\textup{)}.
\item[(b)] $j_1=1$, with $q_*=4.558$ \textup{(}\cref{thm:jone}\textup{)}.
\item[(c)] $j_1=\frac32$, with $q_*=3.7$ \textup{(}\cref{thm:z32}\textup{)}, \zstatus.
\item[(d)] $j_1=2$, with $q_*=4.5579607$ \textup{(}\cref{thm:jtwo}\textup{)}.
\item[(e)] Any $j_1$, if for some pairing of the legs $2$--$5$ the recoupling region has a
  closed colour class and either the two pair dimensions have odd sum or one pair
  contains a spin $\frac12$; then $q_*=2\sqrt\mu\ge4\sqrt{2s(s+1)/\pi}$, with $\mu$ the
  rescaling constant of \cref{lem:mu}
  \textup{(}\cref{thm:closed}\textup{)}. For integer $j_1$ this covers every tuple
  except those whose legs $2$--$5$ are all integers or all half-integers at least
  $\frac32$.
\item[(f)] Any $j_1$, if one of the legs $2$--$5$ has spin $\frac12$; then $q_*=2.7009$,
  and $q_*=4\sqrt2$ if $j_1\ge1$ \textup{(}\cref{cor:ladder}\textup{)}.
\item[(g)] Any $j_1$, if three of the legs $2$--$5$ have spin at most $3$; then
  $q_*=4.5579607$, and moreover $q\ge4\sqrt{10-6\sqrt2}\,j_1-1370$, with a slope that
  cannot be improved \textup{(}\cref{thm:C6}\textup{)}.
\item[(h)] Every tuple with all spins at most $19$, and every tuple with $j_1=\frac52$ and
  all spins at most $35$, with $q_*=4.5579607$ \textup{(}finite certificates,
  \cref{sec:computer}\textup{)}.
\end{enumerate}
\end{theorem}

\begin{center}\small
\begin{tabular}{>{\raggedright\arraybackslash}p{3.6cm}>{\raggedright\arraybackslash}p{3.5cm}>{\raggedright\arraybackslash}p{2.2cm}>{\raggedright\arraybackslash}p{5.5cm}}
\toprule
\estabhead{case} & \estabhead{proved $q_*$} & \estabhead{min.\ found} & \estabhead{method}\\
\midrule
(a) $j_1=\frac12$ & $2.7009$ & $4.5579607$ & analytic; two constants in interval arithmetic\\
(b) $j_1=1$ & $4.558$ & $7.3534$ & analytic; finite part exact\\
(c) $j_1=\frac32$ & $3.7$ & $7.4763$ & analytic; finite box in interval arithmetic\\
(d) $j_1=2$ & $4.5579607$ & $10.5410$ & analytic; finite box in interval arithmetic\\
(e) closed colour class & $4\sqrt{2s(s+1)/\pi}$ & --- & analytic\\
(f) a leg of spin $\frac12$ & $2.7009$; $4\sqrt2$ for $j_1\ge1$ & --- & analytic; windows $5\le N\le149$ in interval arithmetic\\
(g) three legs at most $3$ & $4.5579607$ and $4\sqrt{10-6\sqrt2}\,j_1-1370$ & --- & analytic; exact and interval certificates\\
(h) spins at most $19$; $j_1=\frac52$, spins at most $35$ & $4.5579607$ & $4.5579607$ & rigorous floating point; exact\\
\bottomrule
\end{tabular}
\end{center}

\noindent The column ``min.\ found'' gives the smallest value over all tuples of the slice
with spins at most $19$; it is not a proved infimum.
The smallest value we have found anywhere is $2\sqrt{18-2\sqrt{41}}=4.5579607392$, at
$(\frac12;\frac12,1;1,1)$, where the operator is a four-site chain with squared links
$32,32,80$ and characteristic polynomial $x^4-144x^2+2560$; the minimum at $j_1=\frac52$
over the same range is $12.8129$. So none of the proved constants is the minimum of
its slice; $4.5579607$ is, up to $4\times10^{-8}$, the conjectured global minimum of
\cref{sec:open}, used as a common threshold.

\subsection{The guiding picture: two Stark matrices and a cut}\label{subsec:heronstark}

\Cref{prop:grid} writes every link of $T$ as a Stark link of one pair times
$\frac14\sqrt{F\,G}$, where $F$ depends only on the diagonal $d=b-c$ and $G$ only on
the antidiagonal $u=b+c$. A function of $d$ times a function of $u$ is exactly what a
diagonal congruence (a rescaling of the sites) can produce, so the Heron factors can be
scaled away. Look at
\cref{fig:grid} with this in mind: every horizontal link carries $h_1(b)$, every
vertical link $h_2(c)$, and the factors $\sqrt{F(d\mp\frac12)}$ and $\sqrt{G(u)}$ are
the same along each diagonal and each antidiagonal. Rescale each site by a number that
depends on its two diagonals and these factors disappear; what remains is
$h_1\otimes1+1\otimes h_2$, the Kronecker sum of the two Stark matrices, with the
links that cross the walls or the corner removed. For $s=\frac12$ the strip is the
zigzag chain of \cref{sec:closing}, and the rescaled chain has the Stark links of the
two pairs, alternately.

\begin{theorem}{Heron--Stark congruence}{heron-stark}
For every spin tuple and the pairing $(23|45)$, the entrywise nonnegative matrix $W$ of $T$
satisfies $W=\frac14SW_0S$, where $W_0=P_RKP_R$ is the compression to the region $R$
of the Kronecker sum $K=H(n_1,m_1)\otimes1+1\otimes H(n_2,m_2)$ on the rectangle, and
$S=\diag(\vartheta)$ with $\vartheta(b,c)^2=\phi(d)\,\omega(u)>0$, where
$\phi(d)\phi(d+1)=F(d+\frac12)$ and $\omega(u-1)\omega(u)=G(u)$.
\end{theorem}

\begin{proof}
An $X$-link joins $(b-1,c)$ and $(b,c)$, so it changes $d$ from $d-1$ to $d$ and $u$
from $u-1$ to $u$, and $\vartheta\vartheta'=\sqrt{\phi(d-1)\phi(d)\omega(u-1)\omega(u)}=
\sqrt{F(d-\frac12)G(u)}$; a $D$-link gives $\sqrt{F(d+\frac12)G(u)}$ in the same way.
By \cref{prop:grid} what remains of each link is $\frac14h_1(b)$ or $\frac14h_2(c)$,
the link of $K$. Links of $K$ that leave $R$ are absent from $W_0$ by definition, and
$F$, $G$ are positive inside $R$.
\end{proof}

So $T$ and the compressed Kronecker sum have the same inertia and the same nullity; their
kernels correspond through $\ker W=S^{-1}\ker W_0$. By \cref{blem:starkspec} the uncompressed sum $K$ has spectrum
$\{\lambda+\lambda'\}$ with $\lambda,\lambda'$ running over two Stark spectra: odd
integers if the pair dimensions $N_P=n_P-m_P=2\min(j_a,j_b)+1$ have odd sum, even
integers (with a kernel) if they have even sum. All the difficulty of the
two-dimensional problem is in the cut $P_R$, the triangle inequalities of the
recoupling, and in how much the rescaling $S$ costs.

\begin{lemma}[Scaling with kernels]\label{lem:scalekernel}
For any real matrix $B$ and positive diagonal $S_E$, $S_O$, every nonzero singular value
of $S_EBS_O$ is at least $\min S_E\cdot\min S_O$ times the smallest nonzero singular
value of $B$.
\end{lemma}

\begin{proof}
Put $M=B^TS_E^2B$, which is positive semidefinite. The nonzero eigenvalues of
$S_OMS_O$ are those of $M^{1/2}S_O^2M^{1/2}$. This matrix dominates $(\min S_O)^2M$ and
has the kernel of $M$, so Courant--Fischer on the range compares its nonzero eigenvalues
with those of $M$. Apply this twice.
\end{proof}

\begin{lemma}[The cost of the rescaling]\label{lem:mu}
Let $\mu=\min_{\rm one\ colour}\vartheta^2\cdot\min_{\rm other\ colour}\vartheta^2$. Then
$\mu\ge(2s+2)\Phi_s\ge8s(s+1)/\pi$, where $\Phi_s=2s$ for integer $s$ and
$\Phi_s=2s/\mathcal W_{s-1/2}$ for half-integer $s$, with Wallis's product
$\mathcal W_K=\prod_{l\le K}4l^2/(4l^2-1)<\pi/2$. Hence
\[
  q\ \ge\ 2\sqrt\mu\;\lambda(W_0),
\]
with $\lambda(W_0)$ the smallest nonzero singular value of the bipartite block of $W_0$.
\end{lemma}

\begin{proof}
The free constants of $\phi$ and $\omega$ cancel in $\mu$. Since $G$ increases for
$u>-\frac12$, $\omega(u+2)/\omega(u)=G(u+2)/G(u+1)\ge1$, so the minima over the two
parities of $u$ are $\omega(u_{\min})$ and $\omega(u_{\min}+1)$, with product
$G(u_{\min}+1)\ge2s+2$ because $u_{\min}\ge s$. Along each parity class of $d$,
$\phi(d+2)/\phi(d)=F(d+\frac32)/F(d+\frac12)$ is nonincreasing because $F$ is
log-concave, so $\log\phi$ is concave there and its minimum is at an end; the end values
over $[-s,s]$ are equal for integer $s$ (by $F(x)=F(-x)$), giving $\Phi_s=F(s-\frac12)
=2s$, and differ by Wallis's product $\mathcal W_{s-1/2}$ for half-integer $s$. A
site's colour, the parity of $u-s$, fixes the parity of $u$; and since $u-d=2c$ and
$c$ runs over $m_2+\ZZ$, the parity of $d$ is then fixed too. So $\mu$ is at least the
product. The last statement is \cref{thm:heron-stark,lem:scalekernel}.
\end{proof}

\subsection{Closed regions and the ladder}\label{subsec:closed}

The cut can cost something only where it removes a link of $K$, and it removes links
only at the sites next to the walls $d=\pm s$ and the corner $u=s$; the outer edges of
the rectangle are not cuts, since $K$ itself lives on the rectangle. A colour class
$C$ of the region is \emph{closed} if every neighbour in the rectangle of every site of
$C\cap R$ lies in $R$. The cut sites of the region lie on the wall $d=s$, which has
colour $2(j_4+j_5)\bmod2$; on the wall $d=-s$, of colour $2(j_2+j_3)\bmod2$; and at the
corner $u=s$, of colour $0$ (the colour being the parity of $u-s$). The class of colour
$k$ is closed exactly when no cut site has colour $k$.

\begin{theorem}{Closed regions}{closed}
Suppose a colour class $C$ is closed, and either $N_1+N_2$ is odd or one of the two
pairs contains a spin $\frac12$. Then every nonzero singular value of the bipartite
block of $W_0$ is at least $1$, and every nonzero eigenvalue satisfies
$q\ge2\sqrt\mu\ge4\sqrt{2s(s+1)/\pi}$.
\end{theorem}

\begin{proof}
For $\psi$ supported on $C\cap R$, $K\psi$ is supported on neighbours of $C\cap R$,
all in $R$ by closedness, so the block of $W_0$ from $C$ is $K$ restricted to
$\ell^2(C\cap R)$. If $N_1+N_2$ is odd the spectrum of $K$ consists of odd integers, so
$\|K\psi\|\ge\|\psi\|$. If the first pair is $(\frac12,j)$, its Stark matrix is
$2\times2$ with link $1$, the rectangle is two rails $b=j\pm\frac12$, the two sites
of a position $c$ have opposite colours, and $K$ restricted to $C$ acts as
$\hat\psi\mapsto(H'\pm1)\hat\psi$ on the positions $J$ of $C$, where $H'$ is the
Stark matrix of the other pair. If $N'$ is odd, $H'\pm1$ has odd-integer spectrum. If
$N'$ is even and $J$ is the whole range of $H'$, the nonzero singular values of
$H'\pm1$ are at least $2$. If $J$ misses an end, \cref{lem:onesided} gives
$\|(H'\pm1)\hat\psi\|\ge\|\hat\psi\|$. In every case \cref{lem:mu} finishes.
\end{proof}

The one step in this proof that is not spectral bookkeeping is the case of a chain
with an end removed. Removing an end site from a Stark matrix cannot create a
near-zero mode, because the end components of its eigenvectors are Clebsch--Gordan
weights, largest in the middle of the spectrum and decreasing away from it; so a
vector that avoids the end cannot concentrate on the eigenvalue closest to the shift.

\begin{lemma}[One-sided Stark bound]\label{lem:onesided}
Let $H=H(n,m)$ have even dimension $N$. If $\psi$ vanishes at one end site of the
chain, then $\|(H\pm1)\psi\|\ge\|\psi\|$.
\end{lemma}

The proof is in \cref{app:windows}: $Q=(H-1)^2-1=H(H-2)$ has exactly one negative
eigenvalue, at the eigenvalue $1$ of $H$, and a form with one negative eigenvalue is
nonnegative on the hyperplane $e^\perp$ if and only if $e^TQ^{-1}e\le0$; for the end
vector $e$ this is a sum over the eigenvalues weighted by the squared end components,
and their monotonicity settles the sign.

When no colour class is closed but a leg carries spin $\frac12$, the block is not a
restriction of $H'\pm1$ but a shifted window of it: rows $[a+1,b+1]$ against columns
$[a,b]$. These windows have a uniform bound too.

\begin{theorem}{Shifted Stark windows}{SW}
For every Stark matrix $H(n,m)$ and every window $0\le a\le b\le N-2$, the smallest
singular value of rows $[a+1,b+1]$ times columns $[a,b]$ of $H-1$ is at least
$c_*=\sqrt{(15-\sqrt{97})/6}=0.92657$, with equality at $N=3$, $m=0$.
\end{theorem}

The proof (\cref{app:windows}) reduces every window to the full window $[0,N-2]$, turns
the bound into an explicit inequality between two end-weighted resolvent sums, proves it
for even $N$ and for odd $N\ge151$ by hand, and certifies odd $5\le N\le149$ in interval
arithmetic. A fully analytic constant $\frac2\pi\arcsin\frac34=0.5399$ holds for all
windows.

\begin{corollary}[Ladders]\label{cor:ladder}
If one of the legs $2$--$5$ has spin $\frac12$, then $q\ge6\sqrt2/\pi$ if $j_1=\frac12$,
$q\ge4\sqrt2$ if $j_1$ is an integer, and $q\ge2\sqrt\mu\,c_*\ge6.215$ if $j_1\ge\frac32$
is a half-integer.
\end{corollary}

\begin{proof}
For $j_1=\frac12$ this is \cref{thm:B2}. For integer $j_1$ the number of half-integer
spins among the legs $2$--$5$ is even; pairing them together puts both walls in colour
$0$, the colour of the corner, so the other class is closed, and \cref{thm:closed}
applies with the spin-$\frac12$ pair. For half-integer $j_1\ge\frac32$ either some class is
closed, or (\cref{lem:window} in \cref{app:windows}) the block is a shifted window of
$H'+1$, and \cref{thm:SW,lem:mu} give $q\ge2\sqrt\mu\,c_*$ with $\mu\ge45/4$.
\end{proof}

\subsection{The first slice: \texorpdfstring{$j_1=\frac12$}{j1 = 1/2}}

For $j_1=\frac12$ the region is the chain of \cref{sec:closing}, and the question is
the dimerisation of \cref{fig:ssh} for general spins. We use doubled labels and work
with the links of $8T$ itself, so that the eigenvalues are $q$ directly, with no
factor $8$: let $m_X=2|j_2-j_3|$, $n_X=2(j_2+j_3+1)$, likewise
$m_D,n_D$, and $m=\max(m_X,m_D)$, $n=\min(n_X,n_D)$. The links carry the labels
$l=m+1,\dots,n-1$ in order and alternate in type, and a link of type $P$ of $8T$ has
squared weight $c^2=\frac14(n_P^2-l^2)(l^2-m_P^2)$ (\cref{prop:grid} with $s=\frac12$).

\begin{lemma}[Inner and outer pairs]\label{lem:innerouter}
If the chain has an even number $N=n-m$ of sites, its odd-numbered links all belong to
one pair $O$ with $m_O\le m-1$ and $n_O\ge n+1$, and the other pair owns both $m$ and $n$.
For an even link at label $L=m+2\tau$ followed by an odd link,
\[
  \frac{c_{2\tau}^2}{c_{2\tau+1}^2}\ \le\ \frac{n+L}{n+L+2}\cdot\frac{L-m}{L-m+2}\ \le\ \frac{\tau}{\tau+1}.
\]
If $N$ is odd, the odd links belong to the pair owning $n$ and the even links to the pair
owning $m$.
\end{lemma}

\begin{theorem}{The spin-\texorpdfstring{$\frac12$}{1/2} slice}{B2}
For $j_1=\frac12$ and all other spins, every nonzero eigenvalue satisfies
$q\ge6\sqrt2/\pi>2.70$.
\end{theorem}

The proofs are in \cref{app:chains}. For even $N$ the ratio bound telescopes, so
$E_{ij}^2\le(j/i)/c_{2j-1}^2$, and a lower bound on the odd links turns $\|E\|_F^2$ into a
harmonic double sum. For odd $N$ the ratios increase along the chain; deleting the row
where they cross $1$ splits the chain into two even pieces, one read from each end.
The smallest value found on this slice, over spins at most $19$, is $4.5579607$, so the
constant is not sharp.

\subsection{The second slice: \texorpdfstring{$j_1=1$}{j1 = 1}}

For $s=1$ the strip has three diagonals, $d=-1,0,1$, and the two colour classes are the
middle diagonal and the two walls. So the bipartite block $B$ goes from the diagonal
to the walls, and $B^TB$ lives on the diagonal alone. A diagonal site $(\ell,\ell)$
touches only the four wall sites $(\ell\pm1,\ell)$, $(\ell,\ell\pm1)$, and two
neighbouring diagonal sites share exactly two wall sites; so $B^TB$ is tridiagonal: the
two-dimensional problem at $j_1=1$ is a one-dimensional birth--death problem, which
\cref{prop:pencil1} writes down. Each pair of neighbouring diagonal sites sees its two
shared wall sites through two rows, one weighted by $h_1$ and one by $h_2$. In the
nonnegative gauge of \cref{prop:grid} every link is positive, so every plaquette (four
sites around a unit square) has a positive product of links; we call such a gauge
flux-free. With the sign $(-1)^\ell$ on the diagonal, the two rows then combine into
the square of a difference, $A^+(\Phi_\ell-\Phi_{\ell-1})^2$, plus a nonnegative term.
With a plaquette of negative product the cross terms would cancel instead. The
difference form is a weighted Hardy problem, and Hardy's inequality bounds it.

\begin{proposition}{The pencil at \texorpdfstring{$j_1=1$}{j1 = 1}}{pencil1}
Fix a pairing with pairs $(n_1,m_1)$, $(n_2,m_2)$. Let $\mathrm{lo}=\max(m_1,m_2)$, raised
to $1$ if both are $0$, and $\mathrm{hi}=\min(n_1,n_2)-1$. Put
$P_i(e)=\max(0,(n_i^2-e^2)(e^2-m_i^2))$,
\[
  A^+_e=\frac{(\sqrt{P_1(e)}+\sqrt{P_2(e)})^2}{8e},\qquad
  A^-_e=\frac{(\sqrt{P_1(e)}-\sqrt{P_2(e)})^2}{8e},\qquad \mu_\ell=\frac{2\ell+1}{\ell(\ell+1)}.
\]
Then the nonzero values of $(q/8)^2$ are the nonzero eigenvalues of the pencil
\[
  \mathcal E(\Phi)=\sum_{e=\mathrm{lo}}^{\mathrm{hi}+1}\Bigl[A^+_e(\Phi_e-\Phi_{e-1})^2+A^-_e(\Phi_e+\Phi_{e-1})^2\Bigr]
  \quad\text{against}\quad\sum_\ell\mu_\ell\Phi_\ell^2,
\]
with $\Phi_{\mathrm{lo}-1}=\Phi_{\mathrm{hi}+1}=0$.
\end{proposition}

\begin{proof}
For $s=1$ every wall link has $F(\pm\frac12)=2$. In the gauge $(-1)^\ell$ the two rows
that see $\phi_{\ell-1}$ and $\phi_\ell$ give $\frac{\sqrt2}4(h_1\xi-h_2\upsilon)$ and
$\frac{\sqrt2}4(h_2\xi-h_1\upsilon)$ with $\xi=\sqrt{G(2\ell-1)}\phi_{\ell-1}$, $\upsilon=\sqrt{G(2\ell)}
\phi_\ell$, and $(h_1\xi-h_2\upsilon)^2+(h_2\xi-h_1\upsilon)^2=\frac12[(h_1+h_2)^2(\xi-\upsilon)^2+(h_1-h_2)^2(\xi+\upsilon)^2]$.
Substituting $\phi_\ell=z_\ell\Phi_\ell$ with $z_\ell^2=\mu_\ell$, which solves
$\sqrt{G(2\ell-1)}z_{\ell-1}=\sqrt{G(2\ell)}z_\ell$, and $(4\ell^2-1)h_i^2=P_i$ gives the
pencil. Missing sites and links are exactly the zeros of $P_i$ and of $G(1)$.
\end{proof}

\begin{theorem}{The slice \texorpdfstring{$j_1=1$}{j1 = 1}}{jone}
For $j_1=1$ and all other spins every nonzero eigenvalue satisfies $q\ge4.558$.
\end{theorem}

The proof (\cref{app:jone}) uses the Hardy inequality with constant $4$, a nodal-domain
argument for identical pairs (whose kernel is the constants), and explicit estimates of
the Hardy product when $\max(n_1,n_2)\ge33$; the remaining finitely many pairs of pairs
are certified by exact $LDL^T$ factorisation.

\subsection{Half-integer \texorpdfstring{$j_1$}{j1}: zigzags}

For half-integer $s$ the strip $|d|\le s$ has an even number $2s+1$ of diagonals, and
they pair up: the sites on two neighbouring diagonals $d_k$, $d_k+1$, with
$d_k=-s+2k$, form a single zigzag path, which we call a \emph{zigzag}. Rows of the
bipartite block are the sites with $d+s$ even, columns those with $d+s$ odd. Each
zigzag is then a chain of the kind solved in \cref{sec:closing}, and neighbouring
zigzags are coupled through the links between their diagonals. The picture to keep is
a stack of such chains, glued by couplings whose Gram matrices are a constant multiple
of the neighbouring chain's, up to a diagonal correction.

\begin{proposition}{Zigzag structure}{zigzag}
For half-integer $s$ the bipartite block $B$ is block lower-bidiagonal in the zigzag
index: each diagonal block $D_k$ is a single path alternating between rows and columns,
with imbalance $\#\mathrm{rows}-\#\mathrm{columns}\in\{-1,0,1\}$, and the off-diagonal
block $N_k$ couples zigzag $k$ to zigzag $k-1$. Moreover, for adjacent zigzags,
\[
  N_k^TN_k=\mathsf r_k\,D_{k-1}^TD_{k-1}+\Lambda_k,\qquad \mathsf r_k=\frac{F(d_k-\frac12)}{F(d_k-\frac32)},
\]
with $\Lambda_k$ diagonal.
\end{proposition}

\begin{proof}
A row on $d_k$ has its neighbours on $d_k\pm1$; those on $d_k+1$ are the columns of
zigzag $k$ and those on $d_k-1$ the columns of zigzag $k-1$. On each diagonal the sites
form an interval, which gives the path structure and the imbalance. For the identity,
an off-diagonal entry of $D_{k-1}^TD_{k-1}$ (respectively $N_k^TN_k$) between two
adjacent columns is the product of the two links at their unique shared row of zigzag
$k-1$ (respectively $k$). These two rows, $(b,c)$ on $d$ and $(b+1,c-1)$ on $d+2$, form
a plaquette with the two columns, and by \cref{prop:grid} the two link products contain
the same Stark and $G$ factors and differ only in $F(d+\frac12)$ against $F(d+\frac32)$.
\end{proof}

For $j_1=\frac32$ there are two zigzags, and in a suitable ordered pairing neither has
surplus columns, so $B$ is injective and $\lambda(T)=\smin(B)$. Splitting
$\|B(x,y)\|^2=\|D_0x\|^2+\|Nx+D_1y\|^2$ and giving up a fraction of the first term to
absorb the coupling gives, for every $\tau>0$,
\begin{equation}\label{eq:twochain}
  \smin(B)^2\ \ge\ \min\Bigl(\lambda_{\min}\bigl(D_0^TD_0-\tau N^TN\bigr),\ \tfrac{\tau}{1+\tau}\lambda_{\min}(D_1^TD_1)\Bigr),
\end{equation}
and by \cref{prop:zigzag}, $D_0^TD_0-\tau N^TN=(1-\frac43\tau)D_0^TD_0-\tau\Lambda$. So
everything reduces to two single chains, whose inverses are explicit products of link
ratios, and one diagonal correction. The price is the factor $\tau/(1+\tau)$, paid
once at $j_1=\frac32$. What is left is a bound of $3.777$ at the worst tuple, against a
computed value of $8.87$. The closed-form estimates that make this bound uniform hold only
for large spins (beyond $n_1=4000$ or $10^5$, depending on the regime), so below that
the bound is evaluated parameter by parameter, by computer.

\begin{theorem}{The slice \texorpdfstring{$j_1=\frac32$}{j1 = 3/2}}{z32}
For $j_1=\frac32$ and all other spins every nonzero eigenvalue satisfies $q\ge3.7$.
This is proved \zstatus. Independently of that box, every tuple with $j_1=\frac32$ and
spins at most $20$ satisfies $q\ge7.18$, by exact $LDL^T$ certificates.
\end{theorem}

The proof (\cref{app:z32}) sorts the tuples into eight regimes by the gaps between the
sorted spins of legs $2$--$5$, chooses in each regime a pairing in which both zigzags have
no surplus column, bounds the two chains and the coupling by explicit telescoping
products, evaluates the resulting bound on a finite box of parameters in
outward-rounded interval arithmetic, and controls everything beyond the box by written
tail lemmas. The binding regime is the family $(\frac32;j,j+\frac12,j+\frac12,j+\frac12)$
at $j=1$.

\subsection{The slice \texorpdfstring{$j_1=2$}{j1 = 2}: a triangular reduction}

For integer $j_1$ the strip has an odd number of diagonals, so pairing them into
zigzags leaves one over. Instead, take the colour class whose diagonals are
$d=\pm1,\pm3,\dots$ as columns and let each row diagonal be \emph{owned} by one of its
two neighbouring column diagonals: the owner keeps most of that row's contribution to
$\|Bx\|^2$ and the other neighbour pays a small penalty. Each column diagonal then
carries one tridiagonal problem of the kind solved above. For $j_1=2$ there are only
the two column diagonals $d=\pm1$, the walls $d=\pm2$ are owned by their single
neighbours, and only the middle row diagonal $d=0$ is shared.

Since $T$ links only adjacent diagonals,
\[
  \|Bx\|^2=\sum_\rho\bigl\|Z(\rho-1,\rho)x_{\rho-1}+Z(\rho+1,\rho)x_{\rho+1}\bigr\|^2,
\]
where $Z(\delta,\rho)$ is the zigzag block from column diagonal $\delta$ to row diagonal
$\rho$, a path matrix; and the Gram matrices of the two zigzags at one column diagonal
are proportional off the diagonal, $Z(\delta,\delta+1)^TZ(\delta,\delta+1)-\mathsf r_\delta
Z(\delta,\delta-1)^TZ(\delta,\delta-1)=\diag(\Lambda)$, by the same plaquette argument
as \cref{prop:zigzag}. Applying $\|a+b\|^2\ge(1-\theta)\|a\|^2-(\theta^{-1}-1)\|b\|^2$,
$0<\theta<1$, row diagonal by row diagonal reduces positivity of $B^TB-\lambda$ to
positivity of one tridiagonal matrix per column diagonal.

\begin{theorem}{The slice \texorpdfstring{$j_1=2$}{j1 = 2}}{jtwo}
For $j_1=2$ and all other spins every nonzero eigenvalue satisfies $q\ge4.5579607$.
\end{theorem}

The proof (\cref{app:jtwo}) routes each tuple to a pairing in which the wall chain on the
wall $d=-2$ has no surplus columns, and in one of the three cases neither wall chain does, bounds the chains and the one coupling in closed form, and closes
with an interval-arithmetic box and an explicit tail. Tuples outside the even-sum class
are already covered by \cref{thm:closed}, and tuples with all spins at most $19$ by the
finite certificate.

A proof for all $j_1$ along these lines would need either a uniform version of
\cref{lem:mu} for the compressed Kronecker sum or the large-$j_1$ structure of the next
section; \cref{sec:open} describes how far both go.
